% Removed Sec 4.4 (Projection onto Achievable Centroids), kept for reuse in Sec 5.

\subsection{Projection onto Achievable Centroids}
\label{sec:centroid-projection}

% However, this residual is traded off against the transport cost from the sites, which biases the solution.
% Indeed, at the solution, the volumes are exact and the transport cost equals a constant minus $\smash{2 \sum_i \pvol_i\,\qq_i\tr \centroid_i}$, so that the objective of \eqref{eq:regularized-primal} can be rewritten, up to a constant, as
% 
We now show that the softened problem of \S\ref{sec:regularized-joint} defines a projection map onto the achievable centroids $\achievable$, and characterize when this projection is realized by a power diagram.
% Throughout, we write $\domainCentroid$ also for the configuration with all centroids at $\domainCentroid$.
%
To start, note that, at the solution of $\jointdual_\alpha$, the objective \eqref{eq:regularized-primal-obj} can be rewritten, up to a constant, as
\begin{equation}
  \alpha \sum\nolimits_i \pvol_i \norm{\centroid_i - \bigl( \tcentroid_i + \qq_i / \alpha \bigr)}^2 ,
  \label{eq:regularized-projection}
\end{equation}
which returns the volume-weighted least squares residual between the centroids $\centroid$ and the \emph{pushed targets} $\smash{\tcentroid + \qq/\alpha}$.
Maximizing $\jointdual_\alpha$ thus computes the centroids in $\achievable$ closest to the pushed targets.

Let $\ucentroid$ be centroids verifying \eqref{eq:partition-feasibility-centroid}, which need not be achievable, and let $\qq$ be any sites.
% This entails no loss of generality, since any violation of \eqref{eq:partition-feasibility-centroid} is a uniform translation, which is orthogonal to the affine subspace containing $\achievable$ and thus does not change the closest point.
We define the \emph{projection} of $\ucentroid$ onto $\achievable$ as the closest achievable centroids in the volume-weighted $\smash{L^2}$ sense:
\begin{equation}
  \proj(\ucentroid) = \argmin_{\centroid \in \achievable} \sum\nolimits_i \pvol_i \norm{\centroid_i - \ucentroid_i}^2 .
  \label{eq:projection-map}
\end{equation}
To compute this projection, we maximize $\jointdual_\alpha$ with any sites $\qq$ and prescribed centroids $\tcentroid$ chosen so that the pushed targets coincide with $\ucentroid$, \ie $\smash{\tcentroid = \ucentroid - \qq/\alpha}$.
With these targets, the objective \eqref{eq:regularized-projection} reduces to $\smash{\alpha \sum_i \pvol_i \norm{\centroid_i - \ucentroid_i}^2}$ in which the sites no longer appear and $\alpha$ is only a positive factor, so the optimal partition has centroids $\proj(\ucentroid)$ for any choice of $\qq$ and $\alpha$.
We denote the optimal multipliers of this maximization by $\smash{(\opt{\shift}, \opt{\weight}) = \argmax_{\shift, \weight} \jointdual_\alpha(\shift, \weight)}$, with shifted weighted sites $(\opt{\dual}, \opt{\combweight})$ given by \eqref{eq:shifted-site}, and the projected centroids by $\smash{\opt{\centroid} = \proj(\ucentroid)}$, which are the centroids of the partition $\cell$ generated by $(\opt{\dual}, \opt{\combweight})$.
% 
Since $\achievable$ is convex and compact (\S\ref{sec:background}), this closest point $\opt{\centroid}$ is unique, so $\proj$ inherits the classical properties of projections onto convex sets \cite[Thm.~3.16 and Prop.~4.16]{Bauschke2017Convex}: it is nonexpansive in the volume-weighted norm, and it fixes exactly the achievable centroids, \ie $\smash{\proj(\ucentroid) = \ucentroid}$ if and only if $\smash{\ucentroid \!\in\! \achievable}$.

With exact volumes, the vanishing of the shift gradient \eqref{eq:regularized-gradient-shift} gives $\smash{\shift_i = \alpha\,(\tcentroid_i - \centroid_i)}$, so the shifted sites at the maximizer are
\begin{equation}
  \dual_i
  = \qq_i + \shift_i
  = \qq_i + \alpha \bigl( \ucentroid_i - \tfrac{1}{\alpha}\,\qq_i - \centroid_i \bigr)
  = \alpha\,(\ucentroid_i - \centroid_i) ,
  \label{eq:projection-generators}
\end{equation}
which no longer depend on the sites.
This specific form follows from linearizing the projection objective at its minimizer.
For any partition with volumes $\pvol$, the transport cost to generators $\dual$ expands as $\smash{\int_\domain \norm{\xx}^2 d\xx - 2 \sum_i \pvol_i\,\dual_i\tr \centroid_i + \sum_i \pvol_i \norm{\dual_i}^2}$, so the power diagram of $\dual$ with volumes $\pvol$ selects the achievable centroids that maximize the linear function $\smash{\sum_i \pvol_i\,\dual_i\tr \centroid_i}$ over $\achievable$, \ie $\dual$ is an outward normal of $\achievable$ at its centroids (\S\ref{sec:background}).
On the other hand, the optimality of $\opt{\centroid}$ for the convex problem \eqref{eq:projection-map} states that $\opt{\centroid}$ also minimizes over $\achievable$ the linearization of $\smash{\alpha \sum_i \pvol_i \norm{\centroid_i - \ucentroid_i}^2}$ at $\opt{\centroid}$, which equals $\smash{-2 \sum_i \pvol_i\,\alpha\,(\ucentroid_i - \opt{\centroid}_i)\tr \centroid_i}$ up to a constant, \ie
\begin{equation}
  \sum\nolimits_i \pvol_i\,(\ucentroid_i - \opt{\centroid}_i)\tr (\centroid_i - \opt{\centroid}_i) \le 0
  \quad \text{for all } \centroid \!\in\! \achievable .
  \label{eq:projection-normal}
\end{equation}
Comparing the two linear functions, the generators \eqref{eq:projection-generators} are thus the outward normal $\smash{\ucentroid - \opt{\centroid}}$ of $\achievable$ at the projection, with the magnitude set by the factor $\alpha$ of the penalty, as in \citet[Lemma~4.1]{Bourne2026ConvexOrder}.
While the normal cone alone fixes the generators only up to the homotheties \eqref{eq:lifted-homothety}, the duality selects the representative given by the gradient of the penalty.
As a consequence, the stiffness $\alpha$ only scales the generators uniformly, leaving the power diagram unchanged, while the sites act as a warm start for the maximization, \eg $\smash{\qq = \alpha\,(\ucentroid - \centroid^0)}$ for a guess $\centroid^0$ of the projection.

Although $\proj$ is defined for any centroids $\ucentroid$ verifying \eqref{eq:partition-feasibility-centroid}, its output is realizable only for some of them, and the generators \eqref{eq:projection-generators} distinguish between achievable and unachievable centroids (\S\ref{sec:background}).
If $\ucentroid$ is achievable, then $\smash{\opt{\centroid} = \ucentroid}$ and the generators collapse to $\smash{\dual = 0}$, so the Lagrangian no longer selects a partition, and the optimum is attained by any (generally fractional) partition with centroids $\ucentroid$, which is a power diagram only if $\ucentroid$ is realizable and is never returned as such by the maximization.
If $\ucentroid$ is unachievable, the generators are nonzero, and the optimal partition is the power diagram of \eqref{eq:projection-generators} whenever the normals $\smash{\ucentroid_i - \opt{\centroid}_i}$ are pairwise distinct, in which case $\proj(\ucentroid)$ is realizable, while otherwise the projection lands on a flat face of $\achievable$ exposed by coinciding generators.
Note that the bound \eqref{eq:regularized-distinct} does not apply here, since the targets $\smash{\ucentroid - \qq/\alpha}$ need not lie in the convex hull of $\domain$, so the distinctness of the generators is a property of $\ucentroid$ rather than of the stiffness.
%
In summary, $\proj$ maps unachievable centroids to realizable ones, except when the projection lands on a flat face, and leaves achievable centroids unchanged.

Since $\achievable$ is only described implicitly, deciding whether given centroids are unachievable may seem as hard as the projection itself.
However, any realizable centroids $\centroid$ with generators $\dual$ provide a certificate: since $\centroid$ maximizes $\smash{\sum_i \pvol_i\,\dual_i\tr \centroid_i}$ over $\achievable$, the supporting hyperplane through $\centroid$ with normal $\dual$ separates $\achievable$ from every $\ucentroid$ beyond it, so that
\begin{equation}
  \sum\nolimits_i \pvol_i\,\dual_i\tr (\ucentroid_i - \centroid_i) > 0
  \quad \Longrightarrow \quad
  \ucentroid \!\notin\! \achievable .
  \label{eq:unachievable-test}
\end{equation}
A particular case is the dilation $\smash{\ucentroid = \domainCentroid + \dilation\,(\centroid - \domainCentroid)}$ with $\dilation \!>\! 1$, for which $\smash{\sum_i \pvol_i\,\dual_i\tr (\ucentroid_i - \centroid_i) = (\dilation - 1) \sum_i \pvol_i\,\dual_i\tr (\centroid_i - \domainCentroid)}$.
This quantity is positive since $\domainCentroid$, attained by the uniform fractional partition $\smash{\basis_i = \pvol_i / \domainVol}$, lies in the interior of $\achievable$ relative to \eqref{eq:partition-feasibility-centroid}, while distinct generators are not a uniform translation and thus define a linear function that is not constant on this subspace.
Therefore, dilating realizable centroids about $\domainCentroid$ always yields unachievable centroids, which underlies the projection-dilation method of \citet{Bourne2026ConvexOrder} and is the mechanism we exploit in \S\ref{sec:flow}.

